
All five sub-hypotheses tested in this work pass their respective gates. The core claim is validated.

\subsubsection{5.1 Existence of Hidden State Correctness Signal}

The probe achieves \textbf{AUROC = 0.885} on full-scale evaluation (9,500 train / 1,700 val samples), exceeding the 0.60 random-baseline threshold by a substantial margin.

\subsubsection{5.2 Layer Sweep Results}

The layer sweep confirms the inverted-U pattern:

\begin{table}[htbp]
\centering
\resizebox{\columnwidth}{!}{%
\begin{tabular}{llll}
\toprule
Layer & Depth (\%) & AUROC & 95\% CI \\
\midrule
L3 & 12.5 & 0.629 & [0.53, 0.73] \\
L7 & 25.0 & 0.736 & [0.65, 0.81] \\
L11 & 37.5 & 0.818 & [0.74, 0.88] \\
**L15** & **50.0** & **0.852** & [0.79, 0.90] \\
L18 & 60.0 & 0.822 & [0.75, 0.88] \\
L23 & 75.0 & 0.785 & [0.69, 0.86] \\
L27 & 87.5 & 0.758 & [0.67, 0.83] \\
L31 & 100.0 & 0.766 & [0.67, 0.85] \\
\bottomrule
\end{tabular}
}
\end{table}

Peak performance occurs at L15 (50\% depth). Middle layers (L18 at 60\%) outperform final layers (L31 at 100\%): 0.822 versus 0.766. The full-data probe trained on layer 19 achieves 0.885 AUROC, improving over the layer sweep estimate due to increased training data.

\subsubsection{5.3 Baseline Comparison}

The probe substantially outperforms output-level metrics:

\begin{table}[htbp]
\centering
\resizebox{\columnwidth}{!}{%
\begin{tabular}{lll}
\toprule
Method & AUROC & $\Delta$ vs Probe \\
\midrule
**Probe (L15/L19)** & **0.885** & --- \\
Token Entropy & 0.623 & -0.262 \\
Sequence NLL & 0.589 & -0.296 \\
\bottomrule
\end{tabular}
}
\end{table}

The probe exceeds token entropy by +26.2 AUROC points and sequence NLL by +29.6 AUROC points.

\subsubsection{5.4 Hook Non-Intrusiveness}

Hidden state extraction via forward hooks achieves 100\% output identity (no mismatches across 500 samples) with -3.4\% overhead (hooks were marginally faster than baseline, likely due to caching effects). The extraction is non-intrusive for deployment.

